\documentclass[10pt]{article}
\usepackage[T1]{fontenc}
\usepackage[utf8]{inputenc}
\usepackage[french]{babel}
\usepackage{amsmath,amssymb,mathrsfs,mathpazo}
\usepackage{geometry}
\geometry{a4paper,margin=2cm}
\usepackage{graphicx,array,tabularx,tabularray,colortbl,multirow,multicol,booktabs,enumitem,pifont,ulem,fancybox,diagbox,cancel}
\usepackage{pstricks,pstricks-add,pst-plot,pst-tree,pst-node,pst-text,pst-func,pst-math,pst-bspline,pst-3dplot,pst-eucl,pst-solides3d}
\let\pscalculate\undefined
\usepackage{pst-bezier}
\usepackage{tikz,pgfplots,tkz-tab}
\usetikzlibrary{arrows,arrows.meta,patterns,calc,positioning,trees,decorations.markings,shapes.geometric,shapes.misc,shapes.arrows,intersections,decorations.pathmorphing,decorations.pathreplacing,backgrounds,shadows,babel}
\pgfplotsset{compat=1.18}
\usepgfplotslibrary{fillbetween}
\usepackage[np]{numprint}
\usepackage{listings,makecell,pas-tableur,twoopt}
\let\barre\undefined
\usepackage[pstricks]{tablvar}
\definecolor{darkgreen}{RGB}{0,100,0}
\providecommand{\red}{\color{red}}
\providecommand{\blue}{\color{blue}}
\providecommand{\green}{\color{green}}
\providecommand{\black}{\color{black}}
\providecommand{\white}{\color{white}}
\providecommand{\gray}{\color{gray}}
\providecommand{\cyan}{\color{cyan}}
\providecommand{\magenta}{\color{magenta}}
\providecommand{\yellow}{\color{yellow}}
\providecommand{\violet}{\color{violet}}
\providecommand{\orange}{\color{orange}}
\providecommand{\euro}{\texteuro}
\providecommand{\EUR}{\texteuro}
\providecommand{\ieme}{\textsuperscript{e}}
\providecommand{\degres}{\ensuremath{^\circ}}
\providecommand{\degree}{\ensuremath{^\circ}}
\providecommand{\vv}[1]{\overrightarrow{#1}}
\providecommand{\Coord}[1]{#1}
\providecommand{\decosix}{}
\providecommand{\starredbullet}{\ensuremath{\bullet}}
\pagestyle{empty}
\setlength{\parindent}{0pt}
\usepackage{adjustbox}\DeclareUnicodeCharacter{2236}{:}\DeclareUnicodeCharacter{2190}{\ensuremath{\leftarrow}}
\providecommand{\R}{\mathbb{R}}
\providecommand{\N}{\mathbb{N}}
\providecommand{\Z}{\mathbb{Z}}
\providecommand{\Q}{\mathbb{Q}}
\providecommand{\D}{\mathbb{D}}
\providecommand{\C}{\mathbb{C}}
\providecommand{\e}{{\mathrm{e}}}
\providecommand{\vect}[1]{\overrightarrow{#1}}
\providecommand{\vectt}[1]{\overrightarrow{\mathrm{#1}}}
\providecommand{\vvt}[1]{\overrightarrow{\mathrm{#1}}}
\providecommand{\np}[2][]{{#2~#1}}
\providecommand{\vv}[1]{{\overrightarrow{#1}}}
\renewcommand{\d}{{\mathrm{d}}}
\providecommand{\degres}{{^\circ}}
\providecommand{\degre}{{^\circ}}
\providecommand{\blue}{\color{blue}}
\providecommand{\red}{\color{red}}
\providecommand{\enskip}{\quad}
\providecommand{\llbracket}{\lbrack\!\lbrack}
\providecommand{\rrbracket}{\rbrack\!\rbrack}
\providecommand{\euro}{{\text{€}}}
\providecommand{\no}{{\text{n°}}}
\providecommand{\interval}[2]{{[#1;#2]}}
\providecommand{\intervFO}[2]{{[#1;#2[}}
\providecommand{\dotfill}{\ldots}

\providecommand{\ds}{\displaystyle}
\providecommand{\textding}[1]{#1}
\providecommand{\Oij}{\(\mathrm{O};\vec{\imath},\vec{\jmath}\)}
\providecommand{\Oijk}{\(\mathrm{O};\vec{\imath},\vec{\jmath},\vec{k}\)}
\providecommand{\Ouv}{\(\mathrm{O};\vec{u},\vec{v}\)}
\providecommand{\intervFF}[2]{\ensuremath{[#1;#2]}}
\renewcommand{\euro}{\ifmmode\text{\texteuro}\else\texteuro\fi}
\DeclareUnicodeCharacter{00B0}{\ifmmode^{\circ}\else\textdegree\fi}
\DeclareUnicodeCharacter{20AC}{\ifmmode\text{\texteuro}\else\texteuro\fi}
% Énoncés : Éducation nationale. Transcription et corrigés : APMEP, Denis Vergès et contributeurs. Découpage et classement : CoopMaths, Sébastien Lozano, Éric Elter et contributeurs. Adaptation de la présentation : Maths Explorer.
% CC BY-SA : https://coopmaths.fr/www/about/
% https://coopmaths.fr/alea/static/bac/2023/tex/bac_2023_08_sujet1_caledonie_3.tex
\begin{document}
\noindent\textit{APMEP / CoopMaths. CC BY-SA. Présentation adaptée par Maths Explorer.}\par\medskip
\medskip

\emph{Cet exercice est un questionnaire à choix multiples. Pour chacune des questions suivantes, une seule des quatre réponses proposées est exacte. \\
Pour répondre, indiquer sur la copie le numéro de la question et la lettre de la réponse choisie. \\
Aucune justification n'est demandée. \\
Une réponse fausse, une absence de réponse, ou une réponse multiple, ne rapporte
ni n'enlève de point.}

\medskip

\begin{enumerate}
\item On considère la fonction $f$ définie sur $\R$ par: $f(x) = (x + 1)\text{e}^x$. 

Une primitive $F$ de $f$ sur $\R$ est définie par:
\begin{center}
\begin{tabularx}{\linewidth}{X X}
\textbf{a.~}$F(x) = 1+x\e^x$ &\textbf{b.~}$F(x) = (1+x)\e^x$\\
\textbf{c.~}$F(x) = (2 + x)\e^x$&\textbf{d.~}$F(x) = \left(\dfrac{x^2}{2}  + x\right)\e^x$.
\end{tabularx}
\end{center}
\end{enumerate}
Dans toute la suite de l'exercice, on se place dans l'espace muni d'un repère
orthonormé \Oijk

\begin{enumerate}[resume]
\item On considère les droites $\left(d_1\right)$ et $\left(d_2\right)$ dont des représentations paramétriques sont respectivement:
\begin{center}
$\left(d_1\right)\left\{\begin{array}{l c r}
x&=&2+r\\
y&=&1 + r\\
z&=&- r
\end{array}\right. \:(r \in \R)$\quad et 
$\left(d_2\right)\left\{\begin{array}{l c l}
x&=&\phantom{-}1 - s\\
y&=&-1 + s\\
z&=&\phantom{-}2 - s
\end{array}\right. \:(s \in \R)$\end{center}

Les droites $\left(d_1\right)$ et $\left(d_2\right)$  sont: 
\begin{center}
\begin{tabularx}{\linewidth}{X X}
\textbf{a.~}sécantes.&\textbf{b.~}strictement parallèles.\\
\textbf{c.~}confondues.&\textbf{d.~}non coplanaires.
\end{tabularx}
\end{center}
\item 
On considère le plan $(P)$ dont une équation cartésienne est: 
\[2x - y + z - 1 = 0.\]
On considère la droite $(\Delta)$ dont une représentation paramétrique est:
\[\left\{\begin{array}{l c r}
x&=&2 + u\\
y&=&4 + u\\
z&=&1 - u
\end{array}\right. \:(u \in \R)\]
La droite $(\Delta)$  est:

\begin{center}
\begin{tabularx}{\linewidth}{X X}
\textbf{a.~}sécante et non orthogonale\newline \hspace*{0.3cm} au plan $(P)$.&b. incluse dans le plan $(P)$.\\
\textbf{c.~}strictement parallèle au plan $(P)$.&\textbf{d.~} orthogonale au plan $(P)$.
\end{tabularx}
\end{center}

\item On considère le plan $\left(P_1\right)$ dont une équation cartésienne est $x - 2y + z + 1 = 0$, ainsi que le plan $\left(P_2\right)$ dont une équation cartésienne est $2x + y + z - 6 = 0$.

Les plans $\left(P_1\right)$ et $\left(P_2\right)$ sont:
\begin{center}
\begin{tabularx}{\linewidth}{X X}
\textbf{a.~}sécants et perpendiculaires.&\textbf{b.~} confondus.\\
\textbf{c.~} sécants et non perpendiculaires.& \textbf{d.~} strictement parallèles.
\end{tabularx}
\end{center}

\item On considère les points E(1~;~2~;~1), F(2~;~4~;~3) et G$(-2~;~ 2~;~5)$.

On peut affirmer que la mesure $\alpha$ de l'angle $\widehat{\text{FEG}}$ vérifie:
\begin{center}
\begin{tabularx}{\linewidth}{X X X X}
\textbf{a.~~}$\alpha = 90\degres$&\textbf{b.~\phantom{-}}$\alpha > 90\degres$
 &\textbf{c.~~}$\alpha = 0\degres$&\textbf{d.~~}$\alpha \approx  71\degres$
\end{tabularx}
\end{center}
\end{enumerate}

\bigskip

\end{document}
